\documentclass[11pt,a4paper,english,hyperref]{article}
\usepackage[utf8]{inputenc}

\newcommand\blattNummer{10}
\newcommand\abgabeDatum{10.5.2019}
\newcommand\ersteAufgabeNummer{1}

%%%%%%%%%%%%%%%%%%%%%%%%%%%%%%%%%%%%%%%%%%%%%%%%%%%%%%%%%%%%%%%%%%%%%%%%%% 
\usepackage{babel}
\usepackage{exscale}
\usepackage{amsmath,amsfonts,amssymb}
\usepackage[thref, amsmath, amsthm, thmmarks]{ntheorem}
\usepackage{a4wide}
\usepackage{epic,eepic}
\usepackage{graphicx} % einbinden von Grafiken mit \includegraphics
\usepackage{url}
\usepackage{mathtools}
\usepackage{etoolbox}
\usepackage[colorlinks]{hyperref}
\usepackage{enumitem}
\usepackage[capitalize]{cleveref}
\usepackage{enumitem}

\pagestyle{empty}
\parindent0em
\parskip1ex
\addtolength{\topmargin}{0cm}
\addtolength{\oddsidemargin}{0cm}
\addtolength{\textheight}{2cm}
\addtolength{\textwidth}{0cm}

% neuer Zähler für die Aufgaben
\newcounter{aufgabeNummer}
\setcounter{aufgabeNummer}{\ersteAufgabeNummer}
\addtocounter{aufgabeNummer}{-1}

% in enumerate-Umgebungen erst Kleinbuchstaben, dann kleine römische Zahlen
\renewcommand{\labelenumi}{(\alph{enumi})\,}
\renewcommand{\labelenumii}{(\roman{enumii})\,}
\setenumerate[0]{label=(\alph*)}
\crefname{enumerate}{}{}

% Stil für Definitionen, Sätze usw.
\theoremstyle{break}   % Titelzeile abgesetzt
\newtheorem{Satz}{Proposition}
\newtheorem{Lemma}{Lemma}
\newtheorem{Definition}{Definition}
\theorembodyfont{\upshape}  % normaler Font f"ur Beispiel, Aufgabe
\newtheorem{Beisiep}{Example}
\newtheorem{Aufgabe}[aufgabeNummer]{Exercise}

% Mathematische Kürzel
\newcommand\R{\mathbb{R}}
\newcommand\Q{\mathbb{Q}}
\newcommand\N{\mathbb{N}}
\newcommand\Z{\mathbb{Z}}
\renewcommand\P{\mathbb{P}}
\newcommand\T{{\mathcal T} }

\newcommand\eps{{\varepsilon}}                                  % machine precision?
\newcommand\im{{\sl i\;}}                                        % imaginary unit?
\newcommand\e{{\rm e}}                                          % Euler number
\newcommand\D{{\rm D\;}}
\DeclarePairedDelimiter\abs{\lvert}{\rvert}
\DeclarePairedDelimiter\norm{\lVert}{\rVert}
\DeclarePairedDelimiter\dotProductD{\langle}{\rangle}           % Delimiters for dot product

\DeclareMathOperator{\spann}{span}
\DeclareMathOperator{\cond}{cond}
\DeclareMathOperator{\round}{round}
\DeclareMathOperator{\supp}{supp}

\begin{document}
\parbox{0ex}{ \scalebox{.27}{  \includegraphics[width=.5\textwidth,angle=-90]{cm-logo2.eps} }   }   \\
\parbox{25ex}{
  Prof.~Dr.~S.~Sauter\\
  Institut für Mathematik\\
  Universität Zürich
  }
%
\rule[0cm]{0.cm}{.01cm}                  
\hfill  \parbox{0.6\textwidth}{
  {\sf\LARGE\bfseries Numerik I}
  {\sf\Large\bfseries \;\;--\;\; Homework \blattNummer }\\[1.5ex]
  Deadline: \abgabeDatum,\ 13:00
}
\vspace{5ex}
  
%\textbf{Aufgabe 1} (4 Pkt.)\\
%Die Funktion $f(x)=\sin(x)$ ist nur verrauscht verf\"ugbar: 
%$$
%\tilde f(x) = f(x)+\delta\frac{\sin \frac{1}{x}}{x^2+0.1}
%$$
%
%\begin{enumerate}
%\item Approximieren Sie $f'(0)$ mittels symmetrischer Differenz (d.h.\
%$ f'(x) \approx \frac{\tilde f(x+h)-\tilde f(x-h)}{2h}$).
%\item Wie muss $h$ bez\"uglich $\delta$ optimalerweise gew\"ahlt werden?
%\item Plotten Sie den Fehler der Approximation $e_\delta(h)$ f\"ur $h=10^0,10^{-0.5},\ldots , 10^{-3}$ und f\"ur $\delta=10^{-12},10^{-10},10^{-8},10^{-6}$. \"Uberprufen Sie ob $h=h_{opt}$ aus dem Skript eine gute Wahl ist.
%\end{enumerate}
%\vspace{4ex}
%
%\textbf{Aufgabe 2} (4 Pkt.)\\
%Sei f"ur $m\in \mathbb N$
%$$I_m:=\int_0^1 x^m\ e^x dx.$$
%\begin{itemize}
% \item [(a)] Leiten Sie eine Rekursionsformel zur exakten Berechnung von $I_m$ her und programmieren Sie sie. Testen Sie ihr Programm f"ur $m=20,40,60,80,100$.
%
%\item [(b)] Programmieren Sie die Simpson-Regel und vergleichen Sie die Ergebnisse mit (a). Benutzen Sie dabei 400 Teilintervalle.
%Welche Unterschiede stellen Sie fest und warum? Welches Ergebnis ist f"ur $m=100$ genauer?
%\end{itemize} 
%
%\vspace{4ex}

\begin{Aufgabe}[3 Points, Computational task]
  Given a function $f$, we want to approximate its zeros by a combined bisection and Newton's method. Write a function that takes as input a starting interval and:
  \begin{itemize}
    \item iterates with the bisection method to determine an interval  $I_\delta(a):=[a-\delta, a+\delta]$, where the Newton's method converges to the zero of $f$ for all start values in this interval;
    \item uses Newton's method on $f(x)$ using the starting value $a$ found before.
  \end{enumerate}

  Test your code with the function
  $$f(x)=x-\cos(x)+\sin(x).$$

  \textit{Hint:} For the bisection method, find $\delta$ using Corollary 7.11 from the script. Plot the given function to guess the starting interval to use as the input of your code.
\end{Aufgabe}

\begin{Aufgabe}[8 Points]
  Let $x_*$ be a zero of $f$, i.e. $f(x_*)=0$. We assume:  $f'(x_*)> 0$, $f''(x_*)\neq 0$.
  \begin{enumerate}
    \item (Theoretical task) Find a function $\varphi$ such that $x_*$ is a solution of the fixed-point equation $x=\varphi(x)$, and the iteration formula based on $\varphi$ satisfies
    \begin{equation*}
      x^{(i+1)}=\varphi(x^{(i)}), \quad x^{(i)} \xrightarrow[i\to \infty]{} x_*, 
    \end{equation*}
    with the order of convergence 3. To do this, replace the function $f$ by its Taylor series expansion of order 2 around $x^{(i)}$ and determine $x^{(i+1)}$ as the zero of this parabola. Think about which zero point should be chosen as $x^{(i + 1)}$ if the parabola has two zeros. 
    % Uberlegen Sie sich, welche Nullstelle als $x^{(i+1)}$ gew\"ahlt werden sollte, falls die Parabel zwei Nullstellen hat.
    \item (Computational task) Implement this formula and also the Newton's method to determine $2^{1/3}$. Use the starting value $x_0 = 2$. Compare their convergence rate along with the bisection method using starting interval $[1, 2]$. In all the three methods, approximate up to an accuracy of $10^{-8}$ or maximum of 30 iterations.\\ What are the disadvantages of the method of order 3?
  \end{enumerate}
\end{Aufgabe}

\begin{Aufgabe}[5 Points]
  The function $f(x)=\sin(x)$ is only available in a perturbed (or noisy) form:  
  \begin{equation*}
    \tilde f(x) = f(x)+\delta ~ \dfrac{\sin \frac{1}{x}}{x+1}
  \end{equation*}
  \begin{enumerate}
    \item (Computational task) Approximate $f'(0)$ using symmetric differences (i.e. $ f'(x) \approx \frac{\tilde f(x+h)-\tilde f(x-h)}{2h}$).
    \item (Theoretical task) How should $h$ be choosen optimally in terms of $\delta$?
    \item (Computational task) Plot the error of the approximation $e_\delta(h)$ for $h=10^0,10^{-0.5},\ldots , 10^{-3}$ and for $\delta=10^{-12},10^{-10},10^{-8},10^{-6}$. Check if $h=h_{opt}$ from the script is a good choice.
  \end{enumerate}

\end{Aufgabe}




%\textit{Hinweis:} Benutzen Sie die Taylor-Entwicklung der zweiten Ordnung:
%$$
%f(x)\approx f(x^{(i)})+f'(x^{(i)})(x-x^{(i)})+\frac{1}{2}f''(x^{(i)})(x-x^{(i)})^2.
%$$
%W\"ahlen Sie in jedem Iterationsschritt den Punkt $x^{(i+1)}$, so dass der am n\"achsten zu $x^{(i)}$ steht.

\end{document}